\begin{homeworkProblem}[8]
  Seja $R=C^{\infty}(-1,1)$ o anel das funções $f:(-1,1)\to \mathbb{R}$ infinitamente diferenciáveis, com as operações de soma e produto usual de funções. Mostre que para cada valor de $n$ fixo, o conjunto
  \begin{align*}
    I_{(n)}:=\left\{ f\in R:f^{(k)}(0)=0\text{ para }0\leq k\leq n \right\}
  \end{align*}
  é um ideal de $R$.
  \begin{solution}
  \begin{itemize}
    \item Seja $f,g\in I_{(n)}$, logo note que $(g-f)^{(k)}(0)=g^{(k)}(0)-f^{(k)}(0)=0-0=0$, para todo $0\leq k\leq n$, pelo que podemos conclui que $f-g\in I_{(n)}$.
    \item Seja $f\in I_{(n)}$ e $g\in C^{\infty}(-1,1)$, logo note que $g^{(k)}(0)=c_k\in\mathbb{R}$ para todo $1\leq k\leq n$, então temos que pela regra de Leibniz se cumpre
      \begin{align*}
        (fg)^{k}(0)=\sum_{i=1}^{k}\binom{k}{i}f^{k-i}(0)g^{i}(0)=\sum_{i=1}^{k}\binom{k}{i}f^{k-i}(0)c_i.
      \end{align*}
      Más como para todo $1\leq i\leq k$, temos que $1\leq k-i\leq k\leq n$ o que implica que $f^{{k-i}}(0)=0$, então podemos conclui que
      \begin{align*}
        (fg)^{k}(0)=\sum_{i=1}^{k}\binom{k}{i}f^{k-i}(0)c_i=\sum_{i=1}^{k}\binom{k}{i}(0)c_i=0.
      \end{align*}
      Pelo que se pode conclui que $fg\in I_{(n)}$.\\
      Note que isto é suficiente, ja que $R$ é um anel comutativo, pelo que temos que para cada valor de $n$ fixo, o conjunto $I_{(n)}$ é um ideal de $R$. 
  \end{itemize}
  \end{solution}
\end{homeworkProblem}
